\chapter{State of the art and related work}
\label{chapter:related}

We group related work by the sub-problem it addresses; for each group we state what they solve, how our setting differs
and what we reuse. Table~\ref{tab:related} summarises the comparison.

\section{Policy-as-code and static compliance checking}
Policy engines such as the Open Policy Agent~\cite{opa} express organisational rules in a declarative language (Rego) and
evaluate them against structured inputs (Kubernetes manifests, Terraform plans, API requests); static analysers report
findings in the SARIF interchange format~\cite{sarif2020}. They are deterministic and CI-native, but presuppose that a
human has already translated the regulation into precise rules, and they operate on structured data, not on prose
documentation. Our rule engine is in this tradition (it even evaluates optional Rego checks) and we emit SARIF, but we add
the translation step (requirements derived from clauses) and an LLM for the requirements that no rule can capture.

\section{Requirements extraction and compliance checking with NLP}
Regulatory requirements engineering has long studied how to turn legal text into software requirements: Breaux and
Ant\'on~\cite{breaux2008analyzing} systematically extract rights, obligations and constraints from HIPAA.
Torre et al.~\cite{torre2020completeness} check the completeness of privacy policies against GDPR with NLP and ML, and
Amaral Cejas et al.~\cite{amaral2023dpa} check data processing agreements against GDPR with an NLP pipeline built on a
manually derived set of ``shall'' requirements. Hassani et al.~\cite{hassani2024rethinking} argue that LLMs can replace
much of this hand-crafted pipeline. These works check \emph{documents against documents}; we check a \emph{software
repository} (code, IaC, configuration and documents) against a regulation, and we target the AI Act, which is newer and
has no labelled corpus.

\section{Legal reasoning with LLMs}
LexGLUE~\cite{chalkidis2022lexglue} and LegalBench~\cite{guha2023legalbench} benchmark language models on legal
understanding and reasoning tasks (issue spotting, rule application, interpretation). They show that LLMs perform well on
rule-application tasks but unevenly across task types, which supports using them as \emph{judges of a single, well-scoped
requirement} rather than as end-to-end auditors. Our benchmark is task-specific (requirement $\times$ repository) and
much smaller, but its labels are grounded in repository evidence.

\section{Retrieval-augmented generation and LLM-as-a-judge}
RAG~\cite{lewis2020rag} conditions generation on retrieved passages; we use sparse BM25 retrieval~\cite{robertson2009bm25},
which suits code identifiers and needs no embedding model, and cap the context because long contexts degrade the use of
mid-context information~\cite{liu2024lost}. Zheng et al.~\cite{zheng2023judging} show that strong LLMs agree with human
preferences at a level comparable to inter-human agreement, but exhibit position, verbosity and self-enhancement biases;
we therefore give the judge a structured schema, a single requirement at a time and a hint from the rule engine that it
is explicitly told may be wrong. Calibration of LLM self-reported confidence is imperfect but informative~\cite{kadavath2022know};
we use it only for routing to a human, not as a probability of legal correctness.

\section{Agentic orchestration}
ReAct~\cite{yao2023react} interleaves reasoning with tool calls; AutoGen~\cite{wu2023autogen} composes multiple
conversing agents. Our agents are deliberately \emph{not} free-running: the orchestrator is a fixed program, each agent is
a single structured call without tools (the Claude Code CLI is run with tools disabled), which bounds cost and makes every
step auditable.

\section{Hallucination and citation verification}
Hallucination is pervasive in natural language generation~\cite{ji2023survey}. Attribution can be framed as whether
generated statements are supported by identified sources (AIS)~\cite{rashkin2023ais}; ALCE~\cite{gao2023alce} measures
citation recall and precision of LLM answers; FActScore~\cite{min2023factscore} decomposes outputs into atomic facts checked
against a knowledge source and SelfCheckGPT~\cite{manakul2023selfcheckgpt} detects hallucinations by sampling consistency.
Our Verifier applies the same principle with a cheaper, stricter contract adapted to auditing: citations must refer to
evidence ids actually shown to the judge, positive verdicts require at least one valid citation, and an optional second LLM
pass checks support claim-by-claim.

\section{AI Act scholarship and AI auditing}
Veale and Zuiderveen Borgesius~\cite{veale2021demystifying} analyse the structure and weaknesses of the AI Act proposal,
including its reliance on self-assessment and harmonised standards; Hacker et al.~\cite{hacker2023regulating} discuss
obligations for generative models. Raji et al.~\cite{raji2020closing} propose an end-to-end internal audit framework
(SMACTR) whose artefacts our tool can help produce; M\"okander et al.~\cite{mokander2023auditing} propose layered
(governance, model, application) audits of LLMs. Model cards~\cite{mitchell2019modelcards} and
datasheets~\cite{gebru2021datasheets} are documentation artefacts our evidence collector treats as a separate kind and
our ArtifactGenerator can draft. We position the tool as \emph{support for the internal, application-level audit}, not as a
conformity assessment.

\begin{table}[htbp]
\caption{Comparison with related approaches (\checkmark{} = supported, $\sim$ = partially).}
\label{tab:related}
\centering\small
\resizebox{\linewidth}{!}{%
\begin{tabular}{@{}lcccccc@{}}
\toprule
\textbf{Approach} & \textbf{Legal text} & \textbf{Code/IaC} & \textbf{Open-textured} & \textbf{Citations} & \textbf{Abstains} & \textbf{CI} \\
\midrule
OPA / Rego, SARIF tools~\cite{opa,sarif2020} & -- & \checkmark & -- & \checkmark & -- & \checkmark \\
NLP compliance of documents~\cite{torre2020completeness,amaral2023dpa} & \checkmark & -- & $\sim$ & $\sim$ & -- & -- \\
Legal LLM benchmarks~\cite{guha2023legalbench,chalkidis2022lexglue} & \checkmark & -- & \checkmark & -- & -- & -- \\
LLM + citations~\cite{gao2023alce,min2023factscore} & -- & -- & \checkmark & \checkmark & -- & -- \\
\textbf{Ours} & \checkmark & \checkmark & \checkmark & \checkmark & \checkmark & \checkmark \\
\bottomrule
\end{tabular}}
\end{table}
